% from ega4-1-fr.tex
\oldpage[IV-1]{228}

Pour vérifier qu'un morphisme est quasi-séparé, on est donc ramené à donner
des critères pour qu'un préschéma soit quasi-séparé~:

\begin{proposition}[1.2.7]
\label{IV.1.2.7-fr}
Soient $X$ un préschéma, $(U_\alpha)$ un recouvrement de $X$ formé d'ouverts
quasi-compacts. Les propriétés suivantes sont équivalentes~:
\begin{enumerate}
  \item[a)] $X$ est quasi-séparé.
  \item[b)] Pour tout ouvert quasi-compact $U$ de $X$, l'immersion canonique
  $U\to X$ est quasi-compacte (autrement dit, $U$ est rétrocompact dans $X$).
  \item[b')] L'intersection de deux ouverts quasi-compacts de $X$ est
  quasi-compacte.
  \item[c)] Pour tout couple d'indices $(\alpha,\beta)$, l'intersection
  $U_\alpha\cap U_\beta$ est quasi-compacte.
\end{enumerate}
\end{proposition}

\begin{proof}
Les propriétés \emph{b)} et \emph{b')} sont trivialement équivalentes par
définition d'un morphisme quasi-compact. Comme un ouvert quasi-compact est
réunion finie d'ouverts affines, pour deux ouverts quasi-compacts $U$, $V$ de
$X$, $U\times_{\mathbf Z}V$ est un ouvert quasi-compact de
$X\times_{\mathbf Z}X$ (I, 3.2.7), dont l'image réciproque par $\Delta_X$ est
$U\cap V$, donc \emph{a)} entraîne \emph{b')}. Il est trivial que \emph{b')}
entraîne \emph{c)}~; enfin, si \emph{c)} est vérifiée, les
$U_\alpha\times_{\mathbf Z}U_\beta$ forment un recouvrement de
$X\times_{\mathbf Z}X$ par des ouverts quasi-compacts, et l'on sait que pour
que le morphisme $\Delta_X:X\to X\times_{\mathbf Z}X$ soit quasi-compact, il
suffit que les images réciproques $U_\alpha\cap U_\beta$ des
$U_\alpha\times_{\mathbf Z}U_\beta$ par $\Delta_X$ soient quasi-compactes
(1.1.1)~; donc \emph{c)} entraîne \emph{a)}.
\end{proof}

\begin{corollary}[1.2.8]
\label{IV.1.2.8-fr}
Tout préschéma $X$ dont l'espace sous-jacent est localement noethérien (par
exemple un préschéma localement noethérien) est quasi-séparé~; tout morphisme
$f:X\to Y$ est alors quasi-séparé.
\end{corollary}

\begin{proposition}[1.2.9]
\label{IV.1.2.9-fr}
Soient $X'$, $X''$ deux préschémas, $X=X'\coprod X''$ leur somme,
$f':X'\to Y$, $f'':X''\to Y$ deux morphismes, $f:X\to Y$ le morphisme qui
coïncide avec $f'$ dans $X'$ et avec $f''$ dans $X''$. Pour que $f$ soit
quasi-séparé il faut et il suffit que $f'$ et $f''$ le soient.
\end{proposition}

\begin{proof}
En effet, $X\times_YX$ est somme des quatre préschémas
$X'\times_YX'$, $X'\times_YX''$, $X''\times_YX'$ et $X''\times_YX''$, et
$\Delta_f$ est le morphisme qui coïncide avec $\Delta_{f'}$ dans $X'$ et avec
$\Delta_{f''}$ dans $X''$~; la proposition résulte donc aussitôt des
définitions.
\end{proof}

\subsection{Morphismes localement de type fini}
\label{subsection:IV.1.3-fr}

\begin{env}[1.3.1]
\label{IV.1.3.1-fr}
Rappelons que si $A$ est un anneau, une $A$-algèbre (commutative) $B$ est dite
\emph{de type fini} si elle est engendrée par un nombre fini d'éléments de
$B$, ou, ce qui revient au même, si elle est isomorphe au quotient d'une
algèbre de polynômes $A[T_1,\ldots,T_n]$ par un idéal de cette algèbre. Il est
clair que pour toute $A$-algèbre (commutative) $A'$,
$B\otimes_AA'$ est alors une $A'$-algèbre de type fini. Si $B$ est une
$A$-algèbre de type fini et $C$ une $B$-algèbre de type fini, $C$ est une
$A$-algèbre de type fini, car si $C$ est un quotient de
$B[T_1,\ldots,T_n]$ et $B$ un quotient de $A[S_1,\ldots,S_m]$,
$B[T_1,\ldots,T_n]=B\otimes_AA[T_1,\ldots,T_n]$ est un quotient de
$A[S_1,\ldots,S_m,T_1,\ldots,T_n]$, donc il en est de même de $C$.
\end{env}

\begin{definition}[1.3.2]
\label{IV.1.3.2-fr}
Soient $f:X\to Y$ un morphisme de préschémas, $x$ un point de $X$,
$y=f(x)$. On dit que $f$ est \emph{de type fini en $x$} s'il existe un
voisinage ouvert affine $V$ de $y$ et un voisinage ouvert affine $U$ de $x$
tels que $f(U)\subset V$ et que l'anneau $A(U)$ soit une $A(V)$-algèbre de
type fini. On dit que $f$ est \emph{localement de type fini} s'il est de type
fini en tout point de $X$.
\end{definition}

